% !PRINTED 546
\weberpage{546}
Gehört $\omega$ zum Exponenten $a$, so sind die Potenzen
\weq{9}{1, \omega, \omega^2, \ldots, \omega^{a-1}}
alle incongruent und bilden also, da sie alle der Congruenz (8) genügen (nach 2.), die Gesammtheit der Wurzeln dieser Congruenz. Unter den Zahlen (9) müssen daher alle anderen zum Exponenten $a$ gehörigen Zahlen $\omega$ gesucht werden. Es wird aber $\omega^l$ nur dann zum Exponenten $a$ gehören, wenn $l$ relativ prim zu $a$ ist, und es folgt:

Wenn es überhaupt Zahlen $\omega$ giebt, die zum Exponenten $a$ gehören, so ist ihre Anzahl so gross, wie die Anzahl der relativen Primzahlen zu $a$ in der Reihe der Zahlen $0, 1, 2, \ldots, a-1$. Diese Zahl bezeichnen wir, wie schon früher (Bd. I, \S.\,132), mit $\varphi(a)$. Dass aber zu jedem Theiler $a$ von $p^f-1$ immer wenigstens eine Zahl $\omega$ und folglich $\varphi(a)$ Zahlen gehören, kann man ganz so beweisen, wie der entsprechende Satz der rationalen Zahlentheorie im \S.\,136 des ersten Bandes bewiesen ist.

Die Zahlen $\omega$, die zu dem Exponenten $p^f-1$ gehören, deren es hiernach immer $\varphi(p^f-1)$ giebt, heissen \gesperrt{primitive Wurzeln} von $\pi$. Ist $\gamma$ eine solche primitive Wurzel, so bilden die Potenzen
\[ 1, \gamma, \gamma^2, \ldots, \gamma^{p^f-2} \]
ein volles Restsystem nach dem Modul $\pi$ mit Ausschluss der durch $\pi$ theilbaren Zahl.

Nach dem \gesperrt{Fermat}'schen Lehrsatze für rationale Zahlen ist $t^{p-1}-1$ für $t = 1, 2, \ldots, p-1$ durch $p$, und folglich auch durch $\pi$ theilbar. Die Congruenz
\weq{10}{t^{p-1} \equiv 1 \pmod{\pi}}
hat daher die Wurzeln
\[ 1, 2, \ldots, p-1, \]
und diese sind, da nach \S.\,140,3. eine rationale Zahl nur dann durch $\pi$ theilbar ist, wenn sie durch $p$ theilbar ist, unter einander incongruent. Nach dem Satze 2. hat also die Congruenz (10) keine anderen Wurzeln als diese.

Multipliciren wir die Congruenz (10) noch mit $t$, so folgt, dass die Congruenz $p^{\text{ten}}$ Grades
\[ t^p - t \equiv 0 \pmod{\pi} \]
die $p$ Wurzeln
\[ 0, 1, 2, \ldots, p-1, \]
%CONT>>
